\documentclass[12pt]{article}

% ---------- PAGE SETUP ----------
\usepackage[a4paper,margin=1in]{geometry}
\usepackage{microtype}
\raggedbottom

% ---------- PACKAGES ----------
\usepackage{graphicx}
\graphicspath{{figures/}}
\usepackage{float}
\usepackage{amsmath}
\usepackage{amssymb}
\usepackage{booktabs}
\usepackage{cite}
\usepackage{xcolor}
\usepackage{placeins}
\usepackage[font=small,labelfont=bf]{caption}
\usepackage{titlesec}
\usepackage{hyperref}
\hypersetup{hidelinks}

% ---------- PLACEHOLDER MACRO ----------
% \tbd{...} marks a value that must be replaced with the measured result
% before submission.  Every occurrence is highlighted in red.
\newcommand{\tbd}[1]{\textcolor{red}{\textbf{[#1]}}}

% ---------- SECTION FORMATTING ----------
\titleformat{\subsection}
  {\normalfont\normalsize\bfseries}
  {\thesubsection}{1em}{}
\titleformat{\subsubsection}
  {\normalfont\normalsize\itshape}
  {\thesubsubsection}{1em}{}

% ---------- PARAGRAPH FORMATTING ----------
\setlength{\parskip}{0.5em}
\setlength{\parindent}{0pt}

% ---------- TITLE ----------
\title{A Conversational Multimodal Framework for Early Mental Health Risk
Detection with Dual Elicitation Consistency Modeling and Statistical
Baseline Calibration}
\author{Sengo Basar}
\date{February 2026}

\begin{document}

\maketitle

% =====================================================================
\begin{abstract}
Depression and anxiety frequently remain undetected until they reach an
advanced stage, and existing screening approaches are limited in
complementary ways: clinical instruments are costly and depend on
self-report, while single-modality computational models cannot capture
the distributed nature of mental-health signals. A central difficulty is
the discrepancy between what a person explicitly reports and what their
language, voice, face, and behaviour implicitly reveal.

This paper presents a conversational multimodal framework that treats
that discrepancy as a first-class measurement rather than as noise. The
framework integrates three components. First, a finite state machine
(FSM) interview engine comprising eight clinical modules, grouped into
five thematic domains, elicits explicit ordinal self-reports while
simultaneously capturing implicit linguistic, facial, speech, and
behavioural features. Second, a dual elicitation consistency model
projects implicit representations into the explicit feature space and
quantifies their disagreement as a consistency score $C$. Third, a
statistical baseline layer situates session-level risk and response
latency relative to a population distribution through z-score
normalization, flagging anomalous sessions for human review.

We make a deliberate design decision regarding $C$: although the
empirical results show that high $C$ is associated with lower
classification accuracy, $C$ is not added to the risk score. Instead it
serves as a reliability index and as a gating criterion, so that
uncertainty widens or escalates a case rather than inflating its
apparent severity. A text-level prototype evaluated on a public
depression corpus across three random seeds shows that $C$ is both
label-informative and error-informative, with accuracy declining from
approximately 94\% in the lowest-$C$ tertile to 78\% in the
highest-$C$ tertile. The speech and facial components are specified and
grounded in public benchmarks; their joint evaluation remains future
work. We argue that explicit--implicit inconsistency and
population-relative calibration deserve first-class treatment in
deployable screening systems.
\end{abstract}

\noindent\textbf{Keywords:} Multimodal AI, Mental Health, Speech Analysis,
Facial Emotion, Risk Detection, Consistency Modeling, Dual Elicitation,
Response Bias, Statistical Baseline, Anomaly Detection, Z-Score
Normalization, Valence--Arousal, Dimensional Emotion Modeling, Early Fusion

% =====================================================================
\section{Introduction}

Mental health conditions are difficult to detect because their
presentation is heterogeneous, nonlinear, and frequently masked.
According to the World Health Organization, approximately one in eight
people---roughly 970 million individuals in 2019---lived with a mental
disorder~\cite{who2022wmhr}. More than 720{,}000 people die by suicide
each year, and suicide is the third leading cause of death among
individuals aged 15--29~\cite{who2024suicide}. Depression and anxiety
are also estimated to cost the global economy approximately 12 billion
working days annually~\cite{who2022work}. Because early intervention
substantially improves outcomes, detection at the earliest feasible
stage is a public-health priority.

Current approaches fall short in three distinct ways. First, clinical
instruments such as structured interviews and validated questionnaires
are accurate but expensive, require trained administration, and depend
entirely on what the respondent is willing or able to report. Second,
most computational systems rely on a single modality---text, speech, or
facial expression---even though the relevant signal is distributed
across modalities rather than localised in any one of them. Third, and
most importantly for this work, existing systems do not explicitly
model \emph{disagreement} between signals. When a respondent says ``I am
fine'' while their language becomes more absolutist, their speech slows,
and their facial expressivity flattens, the conflict itself carries
information. Conventional fusion architectures treat such conflict as
noise to be averaged away.

We distinguish two kinds of evidence. \emph{Explicit} evidence is what
the individual consciously and intentionally reports: ordinal responses
to structured questions gathered through a branching conversational
interface. \emph{Implicit} evidence is what the system infers from
behaviour---linguistic markers, response latency, prosody, and facial
dynamics---and is more difficult to deliberately control. We define the
\emph{consistency score} $C$ as the discrepancy between the two
representations. A large $C$ indicates that the individual's explicit
account and implicit behaviour diverge, which is consistent with social
desirability bias, emotional masking, or limited insight.

This paper makes three contributions.
\begin{enumerate}
  \item \textbf{A dual-elicitation conversational architecture.} An
        FSM-based interview engine with eight clinical modules, grouped
        into five thematic domains, that collects explicit self-report
        and implicit multimodal evidence within a single structured
        interaction.
  \item \textbf{Consistency modeling as a first-class signal.} A
        learned projection maps implicit features into the explicit
        feature space, and the residual is quantified as $C$. We
        formalise the \emph{dual role} of $C$---as a label-informative
        signal and as an error-predicting reliability index---and show
        that these roles must not be conflated in a single additive
        risk term.
  \item \textbf{Population-relative statistical calibration.} A
        baseline layer converts absolute risk and latency into
        population-relative z-scores, producing interpretable anomaly
        flags that route ambiguous or atypical sessions to human
        review.
\end{enumerate}

\subsection{Scope of the Claims}
\label{sec:scope}

This is a framework and methods paper accompanied by a \emph{text-level}
prototype evaluation. The prototype uses a public Reddit depression
corpus containing text and behavioural logs; it does not contain speech
or facial recordings. Consequently, the empirical results reported in
Section~\ref{sec:experiments} concern the consistency mechanism and its
statistical properties, not the joint multimodal system. The speech and
facial components are specified in full and grounded in public
benchmarks, but their joint evaluation is explicitly reserved for future
work. We state this limitation up front because earlier drafts of this
work over-claimed multimodal performance on evidence that did not
support it.

% =====================================================================
\section{Related Work}

\subsection{Text-Based Approaches}

Coppersmith et al.~\cite{coppersmith2015mental} demonstrated that
self-disclosed mental-health diagnoses on Twitter can be classified from
language using NLP features and supervised models. Their work established
that linguistic markers carry detectable signal, but it remained
text-only and relied on voluntary public disclosure of diagnosis.
General-purpose contextual encoders such as BERT~\cite{devlin2019bert}
and RoBERTa~\cite{liu2019roberta} substantially improved text
representation learning and are used here as \emph{encoders}, not as
mental-health models in their own right; neither incorporates
non-linguistic signal. DistilBERT~\cite{sanh2019distilbert} retains most
of BERT's task performance at a fraction of the inference cost and is
the encoder adopted in this work.

\subsection{Speech-Based Analysis}

The Distress Analysis Interview Corpus (DAIC-WOZ) was introduced by
Gratch et al.~\cite{gratch2014daic} as a corpus of clinical interviews
with associated depression severity labels. Alhanai et
al.~\cite{alhanai2018daic} subsequently showed that audio and text
sequence models over DAIC-WOZ can detect depression, but the approach
did not integrate facial or behavioural evidence. We cite the corpus and
the modelling work separately to avoid attributing the dataset to the
modelling study.

\subsection{Physiological and EEG-Based Systems}

Picard~\cite{picard1997affective} established affective computing as a
field, arguing that affect should be treated as an input signal amenable
to computational modelling. Koelstra et al.~\cite{koelstra2012deap}
released DEAP, a widely used physiological benchmark for emotion
analysis that has supported SVM, random forest, and CNN baselines.
Physiological sensing nonetheless requires specialised hardware (EEG
headsets, wearable sensors), which limits deployment for low-cost
screening. Russell~\cite{russell1980circumplex} proposed the
circumplex model, representing affect on continuous valence and arousal
axes rather than as discrete categories; this representation underlies
the facial component of the present framework.

\subsection{Facial Emotion Recognition}

Goodfellow et al.~\cite{goodfellow2013fer} released FER2013, an early
large-scale benchmark of 35{,}887 grayscale facial images annotated with
seven emotion classes. Mollahosseini et
al.~\cite{mollahosseini2017affectnet} introduced AffectNet, which added
continuous valence and arousal annotations for in-the-wild images and
therefore supports dimensional rather than purely categorical modelling.
Facial action units, grounded in the Facial Action Coding
System~\cite{ekman1978facs}, provide an alternative descriptive
vocabulary. Kossaifi et al.~\cite{kossaifi2019action} explored action
unit representations on AU-annotated datasets such as BP4D and DISFA,
producing richer descriptors than discrete labels; that line of work
remains facial-only and does not incorporate speech or behaviour.

\subsection{Multimodal Fusion}

Baltru\v{s}aitis et al.~\cite{baltrusaitis2018multimodal} provided the
taxonomy of multimodal machine learning---representation, translation,
alignment, fusion, and co-learning---that organises much of the
subsequent literature. Tensor Fusion Network~\cite{zadeh2017tensor} is a
representative fusion architecture: it models unimodal, bimodal, and
trimodal interactions through an outer product of modality embeddings.
Such architectures assume that modalities are aligned and mutually
reinforcing; disagreement between modalities is treated as noise rather
than as a measurable quantity. MODMA~\cite{liu2022modma} and
related efforts supply clinically collected multimodal data for
depression analysis, while Xu et al.~\cite{xu2025multimodal} fuse BERT
text features with Wav2Vec~2.0 speech features for depression
prediction; their contribution is a better fusion layer, not a model of
agreement or disagreement.

\subsection{Consistency, Response Bias, and Masking}

The discrepancy between self-report and inferred state has a long
history in psychological measurement. Social desirability
bias---the tendency to present oneself favourably---was formalised by
Crowne and Marlowe~\cite{crowne1960social} and remains a documented
source of error in self-report instruments. In the mental-health
context, the same phenomenon appears as emotional masking or
minimisation. Existing computational systems either ignore this
discrepancy or treat it as label noise. We are not aware of a widely
adopted framework that quantifies explicit--implicit inconsistency as
an explicit, inspectable feature and then uses it to modulate
confidence.

\subsection{Research Gap}
\label{sec:gap}

Three gaps motivate this work. (i) \emph{Inconsistency is not modeled.}
Fusion architectures optimise agreement; they do not represent
disagreement as an outcome of interest. (ii) \emph{Global coordination
is absent.} Single-modality and simple fusion models still lack a
mechanism that reconciles evidence across modalities at a session
level. (iii) \emph{Statistical grounding is missing.} Most systems
report an absolute score with no reference to a population distribution,
so a given score cannot be interpreted as typical or atypical for a
comparable individual. Figure~\ref{fig:research_gap} summarises these
gaps against the reviewed literature.

\begin{figure}[H]
    \centering
    \includegraphics[width=\textwidth]{new.png}
    \caption{Research gap analysis across mental-health detection
    approaches. Text-only, speech-only, physiological, and facial
    systems each leave inconsistency modeling and population-relative
    calibration unaddressed; fusion systems optimise agreement rather
    than representing disagreement.}
    \label{fig:research_gap}
\end{figure}

% =====================================================================
\section{Methodology}

\subsection{System Overview}

Figure~\ref{fig:architecture} shows the end-to-end data flow. A user
completes the FSM interview three times. Repeated administration is
deliberate: a single session may reflect a transient state, whereas
patterns that recur across sessions are more likely to be stable and
therefore more trustworthy as evidence. Each session yields a set of
modality-specific feature vectors, a consistency score, and a calibrated
risk estimate.

\begin{figure}[H]
\centering
\includegraphics[width=0.95\linewidth]{dg.png}
\caption{System architecture. The FSM interviewing engine collects
explicit self-reports and implicit multimodal evidence; features are
fused into a joint vector; the consistency score $C$ is computed in
feature space; population-relative z-score calibration situates the
session against a reference distribution; and the reliability gate
routes high-$C$ or statistically anomalous sessions to human review.}
\label{fig:architecture}
\end{figure}

\subsection{Data Sources and Evaluation Cohorts}
\label{sec:data}

The system draws on two qualitatively different kinds of data, and we
keep their roles strictly separate.

\paragraph{Prototype evaluation data.} The empirical study reported in
Section~\ref{sec:experiments} uses a public Reddit depression
corpus~\cite{redditdepression} with associated behavioural consistency
logs. This corpus contains
\textbf{text and behavioural metadata only}. It contains no speech
recordings and no facial video. All quantitative results in this paper
are therefore \emph{text-level} results.

\paragraph{Component grounding benchmarks.} The speech component is
specified against DAIC-WOZ~\cite{gratch2014daic}, which supplies
clinical interview audio and transcripts with depression severity
labels, and MODMA~\cite{liu2022modma}, which supplies clinically
collected voice recordings with emotional-state annotations. The
publicly distributed DAIC-WOZ release provides audio and transcripts
rather than raw facial video; consequently the facial component is
grounded separately in FER2013~\cite{goodfellow2013fer} and
AffectNet~\cite{mollahosseini2017affectnet}. AffectNet provides over
450{,}000 in-the-wild images annotated with both discrete emotion labels
and continuous valence--arousal values, which is what makes dimensional
modelling feasible. These benchmark datasets inform component design and
extractor selection; they do not constitute a joint multimodal
evaluation of this framework.

\paragraph{Behavioural data.} Per-question response latency, session
duration, interaction count, and ordinal self-report values are logged
by the FSM engine itself. These are available in the prototype.

\subsection{Dual Elicitation Design}

\subsubsection{The Interview Engine}

The engine is a finite state machine
\begin{equation}
\pi = (s_0,\; s_1,\; \ldots,\; s_T), \qquad s_T \in F,
\end{equation}
where each transition is conditioned on the preceding response. The
interview comprises \textbf{eight clinical modules} organised into
\textbf{five thematic domains}: physical well-being (sleep, health),
mood and stress (mood, stress), social and self (social, self),
functioning (study), and safety (safety). Branching logic concentrates
follow-up questions where the preceding answer indicates elevated
concern; for example, a response indicating marked social isolation
triggers additional loneliness-focused items. The collection of
modalities is
\begin{equation}
S = S_{\text{text}} \cup S_{\text{speech}} \cup S_{\text{behaviour}}
    \cup S_{\text{facial}} \cup S_{\text{consistency}}.
\end{equation}

\subsubsection{Phase 1: Structured Elicitation (Explicit Features)}

\begin{equation}
\mathbf{x}^{\text{explicit}} = \bigl[\, x_{\text{sleep}},\; x_{\text{mood}},\;
x_{\text{stress}},\; x_{\text{social}},\; x_{\text{self}},\;
x_{\text{study}},\; x_{\text{health}},\; x_{\text{safety}} \,\bigr]^{\top}
\in \mathbb{R}^{m}, \quad m = 8.
\end{equation}

\noindent These values encode what the user \emph{consciously reports} via
normalised ordinal scale responses.

\subsubsection{Phase 2: Conversational Elicitation (Implicit Features)}

\begin{equation}
\mathbf{x}^{\text{implicit}} = \bigl[\, \mathbf{e}_{\text{embed}},\;
x_{\text{sent}},\; x_{\text{self\_ref}},\; x_{\text{dist}},\;
x_{\text{intensity}} \,\bigr]^{\top} \in \mathbb{R}^{k},
\end{equation}

\noindent where $\mathbf{e}_{\text{embed}} \in \mathbb{R}^{768}$ and
$k = 772$. These values encode what the user's language and behaviour
\emph{implicitly reveal}, independent of conscious self-report.

\subsection{Consistency Modeling}
\label{sec:consistency}

\subsubsection{Learned Projection}

A linear map $\phi:\mathbb{R}^{k}\rightarrow\mathbb{R}^{m}$ projects
implicit features into the explicit feature space:
\begin{equation}
\hat{\mathbf{x}}^{\text{explicit}}
= \phi\bigl(\mathbf{x}^{\text{implicit}}\bigr)
= \mathbf{W}\,\mathbf{x}^{\text{implicit}} + \mathbf{b},
\qquad \mathbf{W}\in\mathbb{R}^{m\times k},\; \mathbf{b}\in\mathbb{R}^{m},
\end{equation}
trained by least squares:
\begin{equation}
\mathcal{L}_{\phi}
= \left\| \mathbf{x}^{\text{explicit}}
        - \phi\bigl(\mathbf{x}^{\text{implicit}}\bigr) \right\|_{2}^{2}.
\end{equation}

During the prototype phase, $\phi$ is initialised heuristically. For
example, the mood dimension is estimated as
\begin{equation}
\hat{x}_{\text{mood}}
= w_{1}\,x_{\text{sent}} + w_{2}\,x_{\text{self\_ref}}
+ w_{3}\,x_{\text{dist}} + b .
\end{equation}

\subsubsection{Inconsistency Score}

\begin{equation}
C = \left\| \mathbf{x}^{\text{explicit}}
        - \hat{\mathbf{x}}^{\text{explicit}} \right\|_{2}
= \sqrt{\sum_{i=1}^{m}\bigl(x_i^{\text{explicit}}
        - \hat{x}_i^{\text{explicit}}\bigr)^{2}} .
\end{equation}

\noindent A cosine-based variant that is insensitive to vector magnitude is
\begin{equation}
C_{\cos} = 1
- \frac{\mathbf{x}^{\text{explicit}} \cdot \hat{\mathbf{x}}^{\text{explicit}}}
       {\left\|\mathbf{x}^{\text{explicit}}\right\|_{2}\,
        \left\|\hat{\mathbf{x}}^{\text{explicit}}\right\|_{2}} .
\end{equation}

A large $C$ indicates that linguistic responses contradict numerical
self-reports, which is consistent with emotional masking or social
desirability bias~\cite{crowne1960social}.

\subsubsection{The Dual Role of $C$---and Why It Is Not Added to Risk}
\label{sec:dualrole}

The empirical study (Section~\ref{sec:experiments}) establishes two
facts about $C$: it correlates with the label distribution, and high $C$
predicts classification error. These facts support two different
interpretations, and a naive design conflates them.

If high $C$ is treated as evidence of severity, it should increase the
risk score. If high $C$ is treated as evidence of unreliability, it
should \emph{not} increase the score, because the score is least
trustworthy precisely where $C$ is largest. The second reading is the
one supported by the data: accuracy falls from approximately 94\% in
the lowest-$C$ tertile to 78\% in the highest-$C$ tertile
(Table~\ref{tab:tertiles}). An additive term of the form
$\text{risk} + \gamma C$ would therefore inflate the estimate exactly
where the estimate is weakest. It would also double-count $C$ whenever
$C$ already appears among the fusion features.

We therefore separate the two functions. $C$ is excluded from the point
estimate and is used instead as (i) a \emph{reliability index} reported
alongside the prediction, and (ii) a \emph{gating criterion} that routes
high-$C$ sessions to human review. The reliability index is estimated
empirically as the local accuracy in a neighbourhood
$\mathcal{B}(C)$ of the consistency score:
\begin{equation}
\hat{\rho}(C)
= \frac{1}{|\mathcal{B}(C)|}
  \sum_{j\in\mathcal{B}(C)} \mathbb{1}\bigl[\hat{y}_j = y_j\bigr] .
\end{equation}

\noindent A session is routed to review when
\begin{equation}
\text{review}(i) = \mathbb{1}\bigl[\,C_i > \tau_{C}\,\bigr],
\label{eq:gate}
\end{equation}
where $\tau_C$ is selected on held-out data to trade coverage against
precision. This preserves the clinically important case in which a
respondent reports low distress while implicit signals disagree: such a
session is escalated for human assessment rather than assigned an
artificially high probability by an additive penalty.

\subsection{Feature Extraction}

\subsubsection{Text Features}

\emph{Semantic embedding.} DistilBERT~\cite{sanh2019distilbert} provides
a six-fold inference speed advantage over BERT while retaining most of
its task performance:
\begin{equation}
\mathbf{e}_{\text{embed}} = \text{DistilBERT}(\text{response})[\texttt{CLS}]
\in \mathbb{R}^{768}.
\end{equation}

\emph{Sentiment polarity.} VADER~\cite{hutto2014vader} returns a
compound polarity score:
\begin{equation}
x_{\text{sent}} = \text{VADER}(\text{response}) \in [-1,\,1].
\end{equation}

\emph{Self-reference ratio.} Elevated first-person singular usage is an
established linguistic marker of depressive
rumination~\cite{coppersmith2015mental}:
\begin{equation}
x_{\text{self\_ref}}
= \frac{\text{count(first-person singular pronouns)}}
       {N_{\text{words}}} .
\end{equation}

\emph{Cognitive distortion score.} Absolutist terms such as ``always'',
``never'', ``nothing'', and ``completely'' are detected via a predefined
lexicon:
\begin{equation}
x_{\text{dist}}
= \frac{\text{count(distortion keywords)}}{N_{\text{words}}} .
\end{equation}

\noindent The implicit vector is therefore
$\mathbf{x}^{\text{implicit}}\in\mathbb{R}^{772}$. The distortion and
self-reference terms capture clinical signal that sentiment polarity
alone does not.

\subsubsection{Facial Features via Dimensional Emotion Modeling}

Discrete categorical labels cannot represent continuously graded
affective states, which limits their usefulness for screening. Following
Russell's circumplex model~\cite{russell1980circumplex} as operationalised
in AffectNet~\cite{mollahosseini2017affectnet}, the facial component maps
CNN emotion probabilities onto continuous valence and arousal and
derives a distress index.

\emph{Emotion probability distribution.} A CNN outputs a softmax
distribution over $K=8$ categories:
\begin{equation}
P(\text{emotion}_i) = \frac{e^{z_i}}{\sum_{j=1}^{K} e^{z_j}},
\qquad \sum_{i=1}^{K} P(\text{emotion}_i) = 1,
\end{equation}
where $i \in \{\text{happy, sad, anger, fear, surprise, disgust, neutral,
contempt}\}$.

\emph{Valence.} Valence $V\in[-1,1]$ encodes the positivity of the
expressed state:
\begin{equation}
V = \sum_{i=1}^{K} P(\text{emotion}_i)\,w_V(\text{emotion}_i).
\end{equation}

\emph{Arousal.} Arousal $A\in[-1,1]$ encodes activation intensity:
\begin{equation}
A = \sum_{i=1}^{K} P(\text{emotion}_i)\,w_A(\text{emotion}_i).
\end{equation}

\noindent The weights $w_V$ and $w_A$ are \textbf{heuristic priors}
chosen to respect the circumplex ordering of each category; they are
not regression coefficients estimated from AffectNet's continuous
annotations. The numerical values used in the prototype are:
$w_V(\text{happy})=+1.0$, $w_V(\text{surprise})=+0.4$,
$w_V(\text{neutral})=0.0$, $w_V(\text{contempt})=-0.7$,
$w_V(\text{sad})=-0.7$, $w_V(\text{fear})=-0.8$,
$w_V(\text{anger})=-0.9$, $w_V(\text{disgust})=-0.95$; and
$w_A(\text{fear})=+0.9$, $w_A(\text{anger})=+0.8$,
$w_A(\text{surprise})=+0.7$, $w_A(\text{happy})=+0.6$,
$w_A(\text{contempt})=+0.3$, $w_A(\text{neutral})=0.0$,
$w_A(\text{disgust})=-0.2$, $w_A(\text{sad})=-0.3$. These priors are a
design choice, and Section~\ref{sec:limitations} records them as a
limitation.

\emph{Distress index.} A per-frame distress score combines negative
valence with activation magnitude:
\begin{equation}
D = (1 - V)\,|A| .
\end{equation}

\noindent For fear ($V=-0.8$, $A=0.9$), $D = 1.62$; for anger
($V=-0.9$, $A=0.8$), $D = 1.52$; for sadness ($V=-0.7$, $A=-0.3$),
$D = 0.51$; for happiness ($V=+1.0$, $A=+0.6$) and neutrality
($V=0$, $A=0$), $D = 0$.

This formulation has an important limitation that we state explicitly:
$D$ collapses to zero both for genuinely positive affect and for flat
or blunted affect. It therefore \emph{cannot} by itself distinguish a
calm individual from one exhibiting emotional blunting. We do not claim
otherwise. A candidate discriminator---session-level expressivity
variance, $\sigma_V$ and $\sigma_A$ computed over the session---is
identified in Section~\ref{sec:future} but is not validated here. The
distress index should be read as a measure of negative high-arousal
affect, not as a measure of flatness.

\emph{Temporal smoothing.} Single-frame predictions are noisy; a sliding
window of $N$ frames is maintained:
\begin{equation}
\bar{V} = \frac{1}{N}\sum_{n=1}^{N}V_n, \qquad
\bar{A} = \frac{1}{N}\sum_{n=1}^{N}A_n, \qquad
\bar{D} = \frac{1}{N}\sum_{n=1}^{N}D_n .
\end{equation}

\noindent A default window of $N=30$ frames (approximately one second at
30\,FPS) suppresses single-frame misclassification while preserving
genuine transitions.

\emph{Micro-behavioural features.} MediaPipe
FaceMesh~\cite{lugaresi2019mediapipe} provides facial landmarks from
which the eye aspect ratio (EAR) and mouth aspect ratio (MAR) are
computed:
\begin{equation}
\text{EAR} = \frac{\lVert p_{\text{upper}} - p_{\text{lower}}\rVert_{2}}
                 {\lVert p_{\text{left}} - p_{\text{right}}\rVert_{2}},
\qquad
\text{MAR} = \frac{d_{\text{upper--lower lip}}}
                  {d_{\text{mouth corner distance}}} .
\end{equation}

\emph{Session-level facial vector.}
\begin{equation}
\mathbf{x}^{\text{face}} = \bigl[\,\bar{V},\; \bar{A},\; \bar{D},\;
\text{EAR},\; \text{MAR}\,\bigr]^{\top} \in \mathbb{R}^{5}.
\end{equation}

\subsubsection{Speech Features}

Speech from the recording session is analysed with
librosa~\cite{mcfee2015librosa} and
openSMILE~\cite{eyben2010opensmile} to obtain MFCCs, pitch, jitter,
shimmer, and energy. These prosodic and spectral features enter the
implicit space. Feature selection is informed by the acoustic analyses
reported on DAIC-WOZ~\cite{gratch2014daic,alhanai2018daic} and
MODMA~\cite{liu2022modma}.

\subsubsection{Behavioural Features}

\begin{equation}
\mathbf{x}^{\text{behavior}}
= \bigl[\,\bar{L},\; \sigma^{\text{sess}}_{L},\; D_{s},\; Q\,\bigr]^{\top}
\in \mathbb{R}^{4},
\end{equation}
where
\begin{equation}
\bar{L} = \frac{1}{T}\sum_{t=1}^{T}L_t,
\qquad
\sigma^{\text{sess}}_{L}
= \sqrt{\frac{1}{T}\sum_{t=1}^{T}\bigl(L_t - \bar{L}\bigr)^{2}} ,
\end{equation}
$L_t$ is the response latency for question $t$,
$D_{s}$ is total session duration, and $Q$ is the number of questions
answered. We use $L$ for latency throughout and reserve $\lambda$ for
no other purpose. All features are standardised prior to fusion.

\subsection{Feature Fusion and Risk Estimation}

\subsubsection{Early Fusion}

All modality vectors are concatenated before risk estimation:
\begin{equation}
\mathbf{x}^{\text{fused}}
= \bigl[\,\mathbf{x}^{\text{explicit}},\; \mathbf{x}^{\text{implicit}},\;
\mathbf{x}^{\text{face}},\; \mathbf{x}^{\text{speech}},\;
\mathbf{x}^{\text{behavior}}\,\bigr]^{\top}.
\end{equation}

\noindent Because $C$ is a function of the same features, including it
here would duplicate information and complicate the interpretation of
the gate. It is therefore computed separately and excluded from
$\mathbf{x}^{\text{fused}}$.

\subsubsection{Base Risk Estimate}

\emph{Prototype phase (weighted linear combination).}
\begin{equation}
r = w_{1}\,\bar{x}^{\text{explicit}}
  + w_{2}\,\bar{x}^{\text{behavior}}
  + w_{3}\,\bar{D},
\qquad
\bar{x}^{\text{explicit}} = \frac{1}{m}\sum_{i=1}^{m}x_i^{\text{explicit}},
\label{eq:linear-risk}
\end{equation}
with $w_{1}=0.55$, $w_{2}=0.25$, $w_{3}=0.20$. These weights are
prototype heuristics and are replaced by a learned model in the second
phase.

\emph{Phase 2 (learned risk models).} XGBoost~\cite{chen2016xgboost} is
the primary learner; random forest is the baseline, and LightGBM is used
as an alternative comparison. The learned estimate is
$\hat{r} = f(\mathbf{x}^{\text{fused}})$, and a PCA projection
$\mathbf{z} = \text{PCA}(\mathbf{x}^{\text{fused}})$ is used for
diagnostic visualisation.

\subsubsection{Decision Rule}

The final decision combines the point estimate, the reliability index,
and the statistical baseline (Section~\ref{sec:baseline}):
\begin{equation}
\text{escalate}(i) =
\mathbb{1}\bigl[\,C_i > \tau_C \;\lor\; |Z_{R,i}| > 2
\;\lor\; |Z_{L,i}| > 2 \;\lor\; \text{safety}_i = 1\,\bigr].
\label{eq:decision}
\end{equation}

\noindent The reported output is the pair $\bigl(\hat{r}_i,
\hat{\rho}(C_i)\bigr)$ together with the escalation flag. The point
estimate is never inflated by an inconsistency penalty.

\subsubsection{Planned Extension: Prediction-Level Divergence}
\label{sec:divergence}

The present consistency score operates in \emph{feature space}. A natural
extension assigns each modality its own risk estimate and measures
disagreement in \emph{prediction space}: $D_{ET}$ between explicit and
text-implicit predictions, $D_{EF}$ between explicit and facial
predictions, and $D_{TF}$ between text and facial predictions. This
allows conflict to be localised to specific modalities rather than
reported as a single scalar. It is not implemented here.

\subsection{Statistical Baseline Modeling and Anomaly Detection}
\label{sec:baseline}

\subsubsection{Population Distribution}

For $N$ reference sessions with risk estimates $R_i$ and mean response
latencies $\bar{L}_i$:
\begin{equation}
\mu_R = \frac{1}{N}\sum_{i=1}^{N}R_i,
\qquad
\sigma_R = \sqrt{\frac{1}{N}\sum_{i=1}^{N}\bigl(R_i - \mu_R\bigr)^{2}},
\end{equation}
\begin{equation}
\mu_L = \frac{1}{N}\sum_{i=1}^{N}\bar{L}_i,
\qquad
\sigma^{\text{pop}}_{L}
= \sqrt{\frac{1}{N}\sum_{i=1}^{N}\bigl(\bar{L}_i - \mu_L\bigr)^{2}} .
\end{equation}

\noindent The reference distribution is estimated empirically from the
evaluation cohort. No synthetic population is constructed; this is
recorded as a limitation because a small reference cohort yields
unstable $\mu$ and $\sigma$ estimates.

\subsubsection{Z-Score Normalization}

\begin{equation}
Z_R = \frac{R_i - \mu_R}{\sigma_R},
\qquad
Z_L = \frac{\bar{L}_i - \mu_L}{\sigma^{\text{pop}}_{L}} .
\end{equation}

\noindent Z-scores provide the population-relative context that absolute
scores lack: a given score is interpretable only against what is typical
for comparable sessions.

\subsubsection{Towards Personalised Baselines}

A population baseline cannot distinguish between-subject from
within-subject variation. An individual who is habitually quiet may be
entirely typical for themselves while deviating from the population. The
planned blend is
\begin{equation}
S_i = \kappa_i\,\frac{R_i - \mu_R}{\sigma_R}
    + (1-\kappa_i)\,\frac{R_i - \mu_i}{\sigma_i},
\end{equation}
where $\mu_i$ and $\sigma_i$ are estimated from the individual's own
longitudinal sessions and $\kappa_i \in [0,1]$ controls the blend.
$\kappa_i$ is \emph{not} set in this work; the formulation is included to
fix notation for the planned extension. Note that $\kappa$ is used here
and nowhere else.

\subsubsection{Anomaly Flags}

A session is flagged when a z-score exceeds two standard deviations, or
when consistency is high, or when the safety module is positive
(Eq.~\ref{eq:decision}). Flagged sessions are routed to a clinically
trained reviewer; they are never auto-diagnosed.

% =====================================================================
\section{Experiments and Results}
\label{sec:experiments}

\subsection{Experimental Setup}

The prototype was implemented with scikit-learn, XGBoost, and VADER.
Evaluation used the public Reddit depression corpus with behavioural
consistency logs described in Section~\ref{sec:data} and summarised in
Table~\ref{tab:cohort}. Because this corpus contains only text and
behavioural metadata, \textbf{no speech or facial results are
reported}. Random train--test splits were repeated
across three seeds (42, 12, 7); accuracy, precision, recall, $F_1$, mean
squared error, and $C$ were recorded. Three seeds are sufficient to
describe variability but not to support strong significance claims, so
we report mean $\pm$ standard deviation and treat inferential tests as
exploratory.

\begin{table}[H]
\centering
\caption{Prototype evaluation data and configuration.}
\label{tab:cohort}
\begin{tabular}{@{}ll@{}}
\toprule
Item & Value \\
\midrule
Corpus & Public Reddit depression dataset \\
Modalities available & Text, behavioural metadata \\
Speech / facial data & Not available in this corpus \\
Samples & \tbd{N} \\
Class balance & \tbd{depressed / non-depressed} \\
Random seeds & 42, 12, 7 \\
Primary model & XGBoost~\cite{chen2016xgboost} \\
Baselines & Random forest, LightGBM \\
Text encoder & DistilBERT~\cite{sanh2019distilbert} \\
Sentiment analyser & VADER~\cite{hutto2014vader} \\
\bottomrule
\end{tabular}
\end{table}

\subsection{Metrics}

\emph{Classification.}
\begin{equation}
\text{Acc} = \frac{TP+TN}{TP+TN+FP+FN}, \quad
\text{Prec} = \frac{TP}{TP+FP}, \quad
\text{Rec} = \frac{TP}{TP+FN},
\end{equation}
\begin{equation}
F_{1} = \frac{2\cdot\text{Prec}\cdot\text{Rec}}{\text{Prec}+\text{Rec}} .
\end{equation}

\emph{Regression.}
\begin{equation}
\text{MSE} = \frac{1}{N}\sum_{i=1}^{N}(y_i-\hat{y}_i)^{2}, \qquad
\text{MAE} = \frac{1}{N}\sum_{i=1}^{N}|y_i-\hat{y}_i|, \qquad
R^{2} = 1 - \frac{\sum_i (y_i-\hat{y}_i)^{2}}{\sum_i (y_i-\bar{y})^{2}} .
\end{equation}

\emph{Facial dimensional accuracy.} Consistent with the AffectNet
protocol~\cite{mollahosseini2017affectnet}, valence and arousal
prediction are evaluated with the concordance correlation coefficient:
\begin{equation}
\rho_c =
\frac{2\rho\,\sigma_{\hat{\theta}}\,\sigma_{\theta}}
     {\sigma_{\hat{\theta}}^{2} + \sigma_{\theta}^{2}
      + (\mu_{\hat{\theta}} - \mu_{\theta})^{2}} .
\end{equation}
\noindent This metric applies to the facial component and is reported
only once that component is evaluated; it is not part of the text-level
results below.

\emph{Consistency analysis.} The relationship between consistency and
the available severity labels is quantified by the Spearman rank
correlation, computed on the rank-transformed variables:
\begin{equation}
\rho_s\bigl(C,\,y\bigr)
= \frac{\operatorname{Cov}\bigl(\mathrm{rank}(C),\,\mathrm{rank}(y)\bigr)}
       {\sigma_{\mathrm{rank}(C)}\,\sigma_{\mathrm{rank}(y)}} .
\end{equation}

\subsection{Results}

\subsubsection{Consistency Carries Label and Error Information}

Figure~\ref{fig:consistency_distribution} shows the distribution of $C$
for depressed and non-depressed samples. The distributions are not
equivalent: non-depressed samples concentrate at low $C$, while
depressed samples occupy a wider range with greater mass at higher $C$.
Spearman correlation between $C$ and the label is
$\rho = \tbd{value}$ ($p = \tbd{value}$). This supports the claim that
explicit--implicit inconsistency is associated with depression, while
the overlap between distributions indicates that $C$ is not a
standalone classifier.

\begin{figure}[H]
    \centering
    \includegraphics[width=0.8\textwidth]{consistency_distribution.png}
    \caption{Distribution of the consistency score $C$ for depressed and
    non-depressed samples. Non-depressed samples concentrate at low $C$;
    depressed samples show a heavier upper tail. The overlap indicates
    that $C$ is informative but insufficient on its own as a
    classifier.}
    \label{fig:consistency_distribution}
\end{figure}

\subsubsection{High Consistency Predicts Error}

Table~\ref{tab:tertiles} stratifies performance by $C$ tertile. Accuracy
falls from approximately 94\% in the lowest tertile to 78\% in the
highest, a 16-point gap. This is the central empirical finding of the
paper: $C$ is a usable \emph{reliability} signal. The correct response
is to gate on $C$, not to add it to the risk score
(Section~\ref{sec:dualrole}).

\begin{table}[H]
\centering
\caption{Classification performance stratified by consistency tertile.
Values are percentages; 94.0 and 78.0 are the observed anchors for the
lowest and highest tertiles. Remaining cells must be populated from the
final run.}
\label{tab:tertiles}
\begin{tabular}{@{}lccccc@{}}
\toprule
$C$ tertile & $n$ & Accuracy & Precision & Recall & $F_{1}$ \\
\midrule
Low (most coherent)    & \tbd{n} & 94.0 & \tbd{p} & \tbd{r} & \tbd{f} \\
Middle                 & \tbd{n} & \tbd{a} & \tbd{p} & \tbd{r} & \tbd{f} \\
High (most inconsistent) & \tbd{n} & 78.0 & \tbd{p} & \tbd{r} & \tbd{f} \\
\bottomrule
\end{tabular}
\end{table}

\subsubsection{Cross-Seed Stability}

Table~\ref{tab:seeds} reports mean performance across the three seeds,
and Figure~\ref{fig:failure_boundary} shows the corresponding failure
boundary. Given the small number of runs, we report dispersion rather
than claiming significance; a paired test across seeds is reported as
exploratory only.

\begin{table}[H]
\centering
\caption{Cross-seed stability (mean $\pm$ standard deviation over seeds
42, 12, 7). Inferential tests are exploratory given three runs.}
\label{tab:seeds}
\begin{tabular}{@{}lcccc@{}}
\toprule
Configuration & Accuracy & $F_{1}$ & MSE & $\rho(C,y)$ \\
\midrule
Seed 42 & \tbd{a} & \tbd{f} & \tbd{m} & \tbd{$\rho$} \\
Seed 12 & \tbd{a} & \tbd{f} & \tbd{m} & \tbd{$\rho$} \\
Seed 7  & \tbd{a} & \tbd{f} & \tbd{m} & \tbd{$\rho$} \\
Mean $\pm$ SD & \tbd{a} & \tbd{f} & \tbd{m} & \tbd{$\rho$} \\
\bottomrule
\end{tabular}
\end{table}

\begin{figure}[H]
    \centering
    \includegraphics[width=0.8\textwidth]{failure_boundary_plot.png}
    \caption{Failure boundary as a function of $C$ across seeds 42, 12,
    and 7. The error rate increases monotonically with $C$, supporting
    its use as a reliability index.}
    \label{fig:failure_boundary}
\end{figure}

\subsubsection{Ablation of the Consistency Mechanism}

Table~\ref{tab:ablation} compares four configurations. The final row is
the proposed design: $C$ is used as a gate and excluded from the feature
vector. The ablation isolates the contribution of the consistency
mechanism without confounding it with the additive formulation rejected
in Section~\ref{sec:dualrole}.

\begin{table}[H]
\centering
\caption{Ablation of consistency usage. Configuration~4 is the design
adopted in this paper. Values must be populated from the final run.}
\label{tab:ablation}
\begin{tabular}{@{}llccc@{}}
\toprule
\# & Configuration & Accuracy & $F_{1}$ & Coverage \\
\midrule
1 & Explicit self-report only & \tbd{a} & \tbd{f} & 1.00 \\
2 & Explicit $+$ implicit linguistic & \tbd{a} & \tbd{f} & 1.00 \\
3 & Configuration 2 $+$ behavioural & \tbd{a} & \tbd{f} & 1.00 \\
4 & Configuration 3, $C$ as reliability gate & \tbd{a} & \tbd{f} & \tbd{c} \\
\bottomrule
\end{tabular}
\end{table}

\subsubsection{Filtering Trade-off}

Figures~\ref{fig:path_ab} and~\ref{fig:precision_recall} show the
consequence of excluding high-$C$ sessions entirely. Precision improves
as the exclusion threshold tightens, but coverage falls. This is the expected
precision--coverage trade-off and it confirms that $C$ identifies a
subset of cases that the model handles poorly rather than a subset that
is simply mislabeled.

\begin{figure}[H]
    \centering
    \includegraphics[width=0.8\textwidth]{path_ab_analysis.png}
    \caption{Effect of retaining only low-$C$ sessions. Excluding
    high-$C$ cases defines a high-trust zone with better accuracy, at
    the cost of coverage.}
    \label{fig:path_ab}
\end{figure}

\begin{figure}[H]
    \centering
    \includegraphics[width=0.8\textwidth]{precision_recall_filtering.png}
    \caption{Precision--coverage trade-off when $C$ is used as an
    exclusion filter. Tightening the threshold increases precision while
    reducing the fraction of sessions that receive a prediction.}
    \label{fig:precision_recall}
\end{figure}

\subsubsection{Summary of Findings}

Three findings are supported by the data. (i) $C$ is associated with
the depression label; the distributions differ but overlap. (ii) $C$ is
strongly associated with classification error; accuracy falls by
approximately 16 percentage points between the lowest and highest
tertiles. (iii) Using $C$ to gate predictions yields a higher-trust
subset at reduced coverage. We do \emph{not} report that multimodal
fusion outperforms single-modality models, because the prototype
contains only text and behavioural data.

% =====================================================================
\section{Discussion}

\subsection{Consistency as Reliability, Not Severity}

The central methodological point of this paper is a distinction that is
easy to miss. The data show that $C$ correlates with the label
\emph{and} that $C$ predicts error. These are compatible: in the
high-$C$ region, the signal is both more prevalent and less clean.
Because the model is least accurate there, $C$ must not be folded into
the risk estimate as an additive penalty. Doing so would raise the
score precisely where the score is least trustworthy, and it would
double-count $C$ if $C$ also entered the feature vector.

Treating $C$ as a gate preserves the clinically important case. A
respondent who reports low distress while their language and behaviour
diverge is not assigned a fabricated high probability; instead the
session is flagged for review. This is the appropriate response to
possible masking: escalate uncertainty to a human, do not manufacture
certainty.

\subsection{Dimensional Facial Modeling}

The facial component departs from categorical classification by
representing affect on continuous valence and arousal axes. This is
motivated by the circumplex model~\cite{russell1980circumplex} and by
AffectNet's annotation protocol~\cite{mollahosseini2017affectnet}.

The claim that dimensional modelling distinguishes calm from emotional
blunting is \emph{not} supported by the present formulation. Both
states map near $V \approx 0$, $A \approx 0$, and therefore both yield
$D \approx 0$. What the model captures is negative high-arousal affect
(fear, anger), and this is what the distress index should be understood
to measure. Distinguishing blunting requires an additional measure of
expressivity dynamics, such as session-level $\sigma_V$ and $\sigma_A$,
which is proposed but untested.

Similarly, the weights $w_V$ and $w_A$ are heuristic priors that respect
the circumplex ordering of the eight categories. They are not derived
from AffectNet's continuous annotations and should not be described as
such. Refitting them by regression against the dimensional targets, and
evaluating with the concordance correlation coefficient, is future work.

\subsection{Population-Relative Interpretation}

Absolute risk scores are difficult to interpret in isolation. The
z-score layer converts a session score into a statement of how unusual
it is relative to a reference population, and the anomaly rule routes
statistically atypical sessions to review. The threshold of two
standard deviations is conventional rather than clinically calibrated;
it should be tuned against outcome data before deployment.

\subsection{Repeated Administration}

Requiring three interview administrations is a design choice with a
clear rationale: a single session may reflect a transient state, and
patterns that recur are more likely to be stable. The cost is burden on
the respondent. Whether the stability gain justifies the burden is an
empirical question that this study does not answer.

% =====================================================================
\section{Ethical Considerations}

\subsection{Data Governance}

The prototype reported here uses publicly released datasets. No new
human participants were recruited, and use of each dataset is governed
by the consent and ethics approvals obtained by the original data
providers. Deployment with real users would require independent ethics
approval, informed consent, and a data-protection agreement.

\subsection{Privacy and Scope of Inference}

Sensitive personal data are anonymised and encrypted. The system does
not provide automated clinical diagnoses. The consistency score is an
internal quantity: it is not surfaced to users in raw form, and it is
not presented as a clinical finding. Anomaly flags are directed
exclusively to clinically trained reviewers. Bias monitoring and
fairness evaluation across demographic groups are required before
deployment.

\subsection{Safety Escalation}
\label{sec:safety}

The interview includes a safety module, and the explicit feature vector
contains a safety dimension. The system must therefore be designed on
the assumption that suicidal ideation will be disclosed.

Anomaly flags for \emph{asynchronous} clinical review are not a
sufficient response to acute risk. We specify a separate, immediate
pathway. If the safety module records a positive response, the
screening flow halts, crisis-contact information and local emergency
service numbers are displayed in-session, and the designated on-call
clinician is notified according to a pre-agreed escalation protocol
with a defined response window. The system does not provide emergency
response itself and is not a crisis service; this must be stated to
users before participation, during informed consent, and on every
escalation screen.

We note that a deployed version handling disclosures of this kind
requires clinical governance, an on-call rota, and documented
follow-up, none of which is evaluated in this paper.

\subsection{Transparency}

Model behaviour, the reliability index, and the escalation rule are
documented and auditable. The heuristic status of the valence and
arousal priors is disclosed (Section~\ref{sec:limitations}).

% =====================================================================
\section{Limitations}
\label{sec:limitations}

\emph{Text-only prototype.} The empirical study uses a corpus without
speech or facial recordings. Claims about the joint multimodal system
are architectural, not empirical. The statement that multimodal fusion
outperforms single-modality models is a hypothesis to be tested in
future work, not a result of this paper.

\emph{Statistical power.} Three random seeds and a single corpus limit
inferential strength. The reported tests should be read as exploratory.

\emph{No raw video in DAIC-WOZ.} The publicly distributed DAIC-WOZ
release provides audio and transcripts; it does not provide the raw
facial video that the facial component would require. Facial modeling is
therefore grounded in FER2013 and AffectNet rather than in interview
data.

\emph{Linear projection.} The map $\phi$ is linear and may fail to
capture nonlinear cross-modal relationships.

\emph{Heuristic affect weights.} $w_V$ and $w_A$ are hand-set priors,
not estimated coefficients.

\emph{Distress ambiguity.} $D = (1-V)|A|$ cannot separate calm from
blunted affect; both yield $D \approx 0$.

\emph{Population distribution.} The baseline is estimated from a small
reference cohort, which makes $\mu$ and $\sigma$ unstable. The
$|Z| > 2$ threshold assumes an approximately Gaussian distribution and
has not been clinically calibrated.

\emph{Cross-cultural validity.} The affect priors derive from a
predominantly Western annotation protocol and may not generalise.

\emph{Conflicting evidence handling.} Prediction-level divergence
detection is specified but not implemented; conflict is captured only
partially, in feature space.

\emph{Environmental sensitivity.} Background noise and lighting affect
speech and facial extraction.

\emph{Not a diagnostic device.} The system is a screening aid and does
not replace professional clinical assessment.

% =====================================================================
\section{Future Work}
\label{sec:future}

\begin{enumerate}
  \item \textbf{Joint multimodal evaluation.} Evaluate the full
        text--speech--facial pipeline on a corpus that contains all
        modalities, and test the hypothesis that multimodal fusion
        outperforms single-modality baselines.
  \item \textbf{Prediction-level divergence detection.} Implement
        $D_{ET}$, $D_{EF}$, and $D_{TF}$ as formalised in
        Section~\ref{sec:divergence}, localising conflict to specific
        modality pairs.
  \item \textbf{Nonlinear projection.} Replace the linear $\phi$ with a
        shallow neural network.
  \item \textbf{Learned affect weights.} Refit $w_V$ and $w_A$ by
        regression against AffectNet dimensional annotations and
        evaluate with the concordance correlation coefficient.
  \item \textbf{Expressivity dynamics.} Validate session-level
        $\sigma_V$ and $\sigma_A$ as discriminators of flat versus calm
        affect.
  \item \textbf{Personalised baselines.} Estimate the blend coefficient
        $\kappa_i$ from longitudinal data and compare
        between-subject against within-subject normalisation.
  \item \textbf{Clinical calibration.} Calibrate the z-score threshold
        and the gate threshold $\tau_C$ against PHQ-9~\cite{kroenke2001phq9}
        and GAD-7~\cite{spitzer2006gad7} outcomes using ROC analysis.
  \item \textbf{Uncertainty and attribution.} Add Bayesian uncertainty
        estimation and SHAP-value attribution.
  \item \textbf{Privacy-preserving deployment.} Explore federated
        learning and on-device inference with clinical partners.
\end{enumerate}

% =====================================================================
\section{Conclusion}

This paper presented a conversational multimodal framework for early
mental health risk detection that treats explicit--implicit
inconsistency and population-relative calibration as first-class
concerns. The framework combines an eight-module FSM interview engine
grouped into five thematic domains, a dual elicitation consistency model
that quantifies disagreement between self-report and inferred behaviour,
a dimensional facial component grounded in Russell's
circumplex~\cite{russell1980circumplex} and
AffectNet~\cite{mollahosseini2017affectnet}, and a z-score baseline
layer that situates each session against a population distribution.

The empirical contribution of this paper is a careful characterisation
of the consistency score. On a text-level prototype, $C$ is associated
with the depression label and, more importantly, predicts classification
error: accuracy falls from approximately 94\% in the lowest-$C$ tertile
to 78\% in the highest. This finding has a direct design consequence.
Because high $C$ marks where the model is least reliable, $C$ must not
be added to the risk score; it should be used as a reliability index
and as a gate that routes uncertain cases to human review. The final
decision rule is therefore
\begin{equation}
\text{escalate}(i) =
\mathbb{1}\bigl[\,C_i > \tau_{C} \;\lor\; |Z_{R,i}| > 2
\;\lor\; |Z_{L,i}| > 2 \;\lor\; \text{safety}_i = 1\,\bigr],
\end{equation}
with the reported output being the pair $\bigl(\hat{r}_i,
\hat{\rho}(C_i)\bigr)$ and the escalation flag.

The speech and facial components are specified and grounded in public
benchmarks, but their joint evaluation is future work. By modelling
response bias, emotional masking, the limits of categorical emotion
labels, and the absence of population-relative interpretation as
architectural concerns rather than afterthoughts, the framework provides
a foundation for calibrated, interpretable, and clinically supervised
screening.

% =====================================================================
\bibliographystyle{ieeetr}
\bibliography{references}

\end{document}
